\section{An effective rank bound for an actual Heisenberg extraction}
\label{sec:heisenberg-effective-rank}

\paragraph{Author freeze and shared derivation.}
D-TD64, question Q2ft, is separate from the finite-signature theorem
D-TD63.  The complete author proof, code and first-output readings
pass; complete root and independent whole-proof reviews also pass.
The parent proposed the empty rank interval and
the use of finitely many spherical presentations.  The author supplied
the actual full-period comparison and the complete spherical-kernel
argument.  Sixfolds supplied the complete finite-table, Schreier-word
and relator-area algorithm, independently whole-checked by the author.
Its full argument is incorporated below, with this shared provenance.
The algorithm is terminating; no claim of a practical running time or
an evaluation of its enormous final constants is made.

\paragraph{The actual-source hypotheses.}
Let a connected finite \(J\)-Galois cover of \(\mathbf P^1\) have a
faithful irreducible complex coefficient representation \(V\), of
dimension \(e\), and its actual generating product-one inertia
operators \(T_1,\ldots,T_r\).  Suppose an integer \(M\ge1\) satisfies
\[
 T_i\ne I,\qquad r\ge12,\qquad
 \sum_{i=1}^r\operatorname{rank}(T_i-I)=2e+6,\qquad
 \operatorname{rank}(T_i-I)\ge e/M\quad(1\le i\le r).
 \tag{HER.1}
\]
Suppose \(V=W\otimes U\) is preserved factorwise and its actual finite
projective image \(P\) on \(W\) is primitive.  It contains a normal
nondegenerate projective Heisenberg subgroup
\[
 Q\simeq\mathbf F_p^{\,2m},\qquad m\ge1,\qquad
 a=\dim W=p^m.
 \tag{HER.2}
\]
The represented action of \(Q\) on \(W\) is irreducible.  The entire
complement \(U\) is retained and is otherwise arbitrary.
The actual label image \(H\subset\operatorname{Sp}(Q)\) is assumed
irreducible on the original \(\mathbf F_p\)-space \(Q\).
This last hypothesis is separate from primitivity of \(P\); it holds
for a minimal-normal Heisenberg extraction in TD12.
No source budget is assigned separately to \(W\) or \(U\).

Retain \(B,C,O\) from accepted TD63 and define the further exact integers:
\[
 \begin{split}
 \ell&=\min\{j\ge0:2^j\ge8M\},\qquad B=\max\{7743,M\},\\
 C&=\max\{64M^2,318\ell,2544\},\qquad O=B^{C^2+1},\\
 A_0&=\max\{24,6M\},\qquad
 L_0=\min\{j\ge0:2^j\ge A_0\},\\
 R&=\min\{j\ge0:2^j\ge4M\},\qquad
 S=\max\{2O,60\},\qquad T_0=2^{8O}R.
 \end{split}
 \tag{HER.3}
\]
Here and below the minima range over integers.

For each ordered tuple \(\boldsymbol d=(d_1,\ldots,d_k)\) with
\(3\le k\le8\), \(1\le d_i<O\), and
\(\sum_i(1-1/d_i)<2\), form the presentation
\[
 \Gamma_{\boldsymbol d}=
 \langle x_1,\ldots,x_k\mid
 x_1\cdots x_k=1,\ x_i^{d_i}=1\ (1\le i\le k)\rangle.
 \tag{HER.4}
\]
Steps 3--4 below give a completely specified terminating algorithm
which, from this tuple and \(S\), produces a positive integer
\(\mathcal D_{\boldsymbol d}\).  It is the maximum of one and the sum
of certified power-relator areas of an ordinary Schreier generating
list for the free presentation kernel.  Set
\[
 \mathcal D_*(O)=
 \max_{\substack{3\le k\le8,\ 1\le d_i<O\\
                  \sum_i(1-1/d_i)<2}}
                   \mathcal D_{\boldsymbol d},
 \qquad
 K(M)=\max\{2L_0,S,\mathcal D_*(O)T_0\}.
 \tag{HER.5}
\]
This is a finite nonempty family.  The displayed definition, together
with the algorithm, defines an integer-valued computable function of
\(M\); it is not an unspecified existence constant.

\begin{theorem}[D-TD64]
Under precisely the actual-source hypotheses above,
\[
             p\le B,\qquad 2m\le K(M),\qquad
             |P|<B^{K(M)^2+K(M)}.
 \tag{HER.6}
\]
Consequently the selected actual projective Heisenberg factor has
bounded prime, rank, dimension and group order in terms of the whole
original support floor \(M\).  Together with TD63, its projective
groups, representations and marked tuples lie in finite effective
enumerations, without any assertion that an enumerated object occurs
in a geometric source.

This theorem does not bound an arbitrary original inducing index,
the original scalar periods, the full original group, or other
uncontrolled factors of \(V\).  It uses irreducibility of the label
module over \(\mathbf F_p\), not absolute irreducibility.
\end{theorem}

\begin{proof}
\emph{Step 1: the finite set of actual projective periods.}
All hypotheses of TD63 hold.  For each original index let \(y_i\in P\)
be the selected projective inertia and let \(g_i\in H\) be its label.
TD63 gives \(p\le B\) and \(o_i=\operatorname{ord}(y_i)<O\), including
the defining-characteristic part of the order.
The kernel of \(P\to H\) is exactly \(Q\), by TD05; a nonidentity
\(y_i\) may therefore have \(g_i=I\).

If \(a\le A_0\), then \(2^m\le p^m=a\le A_0\le2^{L_0}\),
so \(2m\le2L_0\), as required.  Assume for the rest of the rank proof
that \(a>A_0\).  Then \(e\ge a>24\) and \(6/e<1/M\).
Let \(k\) be the number of original indices with \(y_i\ne1\).
The all-parameter quarter floor TD08, transferred through the entire
partner by spectral convolution, implies
\[
       k/4\le 2+6/e<9/4,\qquad 3\le k\le8.
 \tag{HER.7}
\]
The lower bound is a generation statement: a generating product-one
tuple with at most two nonidentity entries generates a cyclic group,
whereas \(P\) contains the noncyclic elementary abelian group \(Q\).
There are at least \(r-k\ge4\) omitted original indices.  Their selected
projective image is identity but their actual coefficient operator is
still genuine, so each retains the cost \(1/M\) in (HER.1).

Relabel the \(k\) nonidentity selected indices in their original order.
They still have product one and generate \(P\), and their labels
generate \(H\).  Here \(2\le o_i<O\) and \(g_i^{o_i}=I\).
Consider all the integers
\(\operatorname{rank}_{\mathbf F_p}(g_i^j-I)\),
\(0\le j<o_i\), for these \(k\) indices.  There are fewer than \(8O\)
of them, counted with repetitions.  The \(8O+1\) disjoint intervals
\((2^jR,2^{j+1}R]\), \(0\le j\le8O\), cannot all contain a rank.
Choose the first empty interval and write \(T=2^jR\).  Thus
\[
 R\le T\le T_0,\qquad
 \operatorname{rank}(g_i^u-I)\notin(T,2T]
          \quad\text{for every }i,u.
 \tag{HER.8}
\]

\emph{Step 2: low-rank powers form subgroups and force a spherical sum.}
For each \(i\), let
\[
 E_i=\{u\in\mathbf Z/o_i\mathbf Z:
              \operatorname{rank}_{\mathbf F_p}(g_i^u-I)\le T\}.
 \tag{HER.9}
\]
It contains zero and is closed under negatives, since an invertible
operator and its inverse have equal moving rank.  If \(u,v\in E_i\),
the rank product inequality gives
\(\operatorname{rank}(g_i^{u+v}-I)\le2T\).  The empty interval in
(HER.8) makes this rank at most \(T\).  Hence \(E_i\) is a subgroup.
Let \(d_i=[\mathbf Z/o_i\mathbf Z:E_i]\).  Then \(d_i\mid o_i\),
\(E_i\) consists exactly of the multiples of \(d_i\), and
\[
   1\le d_i<O,\qquad
   \operatorname{rank}_{\mathbf F_p}(g_i^{d_i}-I)\le T.
 \tag{HER.10}
\]
The case \(d_i=1\), including an identity label, is retained.

Let \(L_i=\operatorname{ord}(T_i)\) be the full original coefficient
order at this index.  Its image in \(P\) has order \(o_i\), so
\(o_i\mid L_i\).  Both actual finite projective factor images have
finite determinant-one lifts, obtained by taking their inverse images
in the corresponding special linear groups.  Average positive
Hermitian forms over these finite lifts.  Representatives from them
are unitary, and the scalar needed to express \(T_i\) as their tensor
product is a root of unity, since both the tensor and \(T_i\) have
finite order.  Absorb that phase in the first factor and write
\(T_i=A_i\otimes B_i\) with unitary factors.  Since
\(T_i^{L_i}=I\), their \(L_i\)-th powers are inverse unit scalars.
Rescale by inverse unit phases to give both factors \(L_i\)-th power
identity.  These phase rescalings change no normalized absolute trace
or projective support.
Thus this Fourier argument uses the full original period, including
every selected-projectively-identity power.

TD05's exact Weyl trace identity gives
\[
 \frac{|\operatorname{tr}(A_i^u)|}{a}
 \le p^{-\operatorname{rank}_{\mathbf F_p}(g_i^u-I)/2}.
 \tag{HER.11}
\]
It holds also for \(p=2\) and for label-identity powers; no conclusion
that such a power is scalar is needed.  Outside \(E_i\), the right
side is strictly less than \(p^{-T}\le2^{-R}\).  On \(E_i\) we use only
the bound one.  In a complete original period, precisely \(1/d_i\)
of the exponents lie in \(E_i\).

For every eigenvalue target of \(T_i\), Fourier inversion and the
unitary bound \(|\operatorname{tr}(B_i^u)|\le\dim U\) therefore give
\[
 \frac{\nu_V(T_i)}e
 \ge (1-1/d_i)(1-c),\qquad c=2^{-R}\le\frac1{4M}.
 \tag{HER.12}
\]
In particular this lower bounds its actual moving rank at target one.
Write \(\sigma=\sum_{i=1}^k(1-1/d_i)\).  The original omitted entries
and (HER.1) imply
\[
  2+\frac6e\ge\sigma(1-c)+\frac4M.
 \tag{HER.13}
\]
If \(\sigma\ge2\), the right side is at least
\(2-2c+4/M\ge2+7/(2M)\).  This contradicts \(6/e<1/M\).
Consequently \(\sigma<2\), and the tuple \(\boldsymbol d\) belongs
to the finite family in (HER.5).

\emph{Step 3: finiteness of the complete spherical presentation.}
We prove that (HER.4) has order at most \(S\), not merely that it has
some finite rotation-group quotient.  Delete each generator with
\(d_i=1\), using its relation \(x_i=1\).  If \(q\) orders remain,
each is at least two and contributes at least \(1/2\) to \(\sigma\),
so \(q\le3\).
For \(q=0\) or1 the group is trivial.  For \(q=2\), product one makes
it cyclic of order dividing the greatest common divisor of the two
orders, hence of order less than \(O\).

The order of the remaining entries can be permuted without changing
the group: the adjacent Hurwitz replacement
\((x,y)\mapsto(xyx^{-1},x)\) preserves the product and exchanges the
two power orders, and has an inverse of the same kind.  For \(q=3\),
sort the orders \(2\le b_1\le b_2\le b_3\).  The inequality is
\(1/b_1+1/b_2+1/b_3>1\).  If \(b_1\ge3\) it fails; thus \(b_1=2\).
If \(b_2\ge4\) it fails, so \(b_2=2\) or3.  The complete list is
\[
 (2,2,d)\ (2\le d<O),\qquad
 (2,3,3),\qquad(2,3,4),\qquad(2,3,5).
 \tag{HER.14}
\]
In the first case eliminate the third generator.  With \(z=xy\),
the relations are \(x^2=1\), \(z^d=1\), and \(xzx=z^{-1}\).
Every word is therefore one of \(z^j\) or \(xz^j\),
\(0\le j<d\); the group has order at most \(2d<2O\).

For completeness, the other three groups have orders 12,24,60 by
the following elementary presentation argument.  Take a regular
tetrahedron, octahedron or icosahedron, respectively, inscribed in
the sphere.  Subdivide each triangular face by its face centre,
edge midpoints and vertices, and project radially onto the sphere.
Reflection across each side of a resulting flag triangle preserves
the corresponding regular polyhedron.  Adjacent flag triangles
are therefore related by a reflection.  The connected adjacency
graph shows transitivity on all flag triangles, while an isometry
fixing a flag fixes three linearly independent rays and is identity.
There are \(6F=24,48,120\) flags for \(F=4,8,20\) faces.  The
orientation-preserving subgroup consequently has order 12,24,60
and two triangle orbits.  Its quotient sphere is explicitly the
double of one closed flag triangle: the two orientations are the
two faces and the three matching sides are glued.

The only nontrivial point stabilizers occur at the three types of
vertices of this quotient triangulation.  Their orders are two at
edge midpoints, three at face centres, and respectively three,
four or five at original vertices.  Interior points of a flag
triangle have trivial stabilizer; side interiors can be fixed only
by the orientation-reversing reflection, not by a nonidentity
rotation.  Thus the quotient map \(S^2\to S^2\) has precisely the
three branch orders in the corresponding case of (HER.14).

Remove those three branch values and all their preimages.  This is
a connected regular unbranched covering of punctured spheres by
the rotation group.  The path-lifting description of the kernel
and deck group is the standard covering-space statement in
\cite[\S1.3, Propositions 1.31--1.32 and 1.39--1.40,
pp.~61 and 71--72]{TD:Hatcher2002}.
Here its presentation consequence can be checked directly.
Choose positively oriented base puncture loops \(x_1,x_2,x_3\)
with product one.  A small loop about an upstairs puncture maps
to a conjugate of \(x_i^{b_i}\), because its local covering degree
is \(b_i\).  The fundamental group of the upstairs punctured
sphere is generated by these peripheral loops: cutting along
arcs to the punctures gives the usual free group with the single
product relation.  Hence the full covering kernel is contained
in the normal closure of the three power relators.  Conversely,
each power relator lifts closed, so that normal closure is
contained in the covering kernel.  They are equal.  Therefore
the full group in (HER.4), not just an image of it, is the indicated
rotation group.  This proves in every case that
\[
                     |\Gamma_{\boldsymbol d}|\le S.
 \tag{HER.15}
\]

\emph{Step 4: an effective finite-presentation rank lemma.}
We give the entire shared relator argument to specify the function
in (HER.5).  This step works over every field and for any irreducible
module.  Let \(F\) be the free group on \(x_1,\ldots,x_{k-1}\), set
\(x_k=(x_1\cdots x_{k-1})^{-1}\), and let \(\mathcal R\) be the
normal closure in \(F\) of the \(k\) words \(q_i=x_i^{d_i}\).
Thus \(F/\mathcal R=\Gamma_{\boldsymbol d}\).

First enumerate all multiplication tables on sets of size at most
\(S\), test all group axioms, and enumerate all generating marked
\((k-1)\)-tuples satisfying the \(k\) relators.  This is a finite
search.  Each surviving marked table is a quotient of
\(\Gamma_{\boldsymbol d}\).  By (HER.15), the full group itself occurs.
Choose a survivor of maximum order, breaking ties lexicographically
in a fixed encoding.  Its order \(u\) equals
\(|\Gamma_{\boldsymbol d}|\), and its marked quotient map is an
isomorphism.  This selection does not assume a word-problem oracle.

In this table choose a breadth-first representative \(t_\gamma\in F\)
for each \(\gamma\), with \(t_1=1\), using a fixed ordering of the
\(2(k-1)\) generator letters and their inverses.  Each representative
has length at most \(u-1\).  Form the at most \(2(k-1)u\) Schreier
words
\[
 w_{\gamma,z}=t_\gamma z\,t_{\gamma\bar z}^{-1},
 \qquad \gamma\in\Gamma_{\boldsymbol d},\quad
 z\in\{x_1^{\pm1},\ldots,x_{k-1}^{\pm1}\}.
 \tag{HER.16}
\]
These lie in \(\mathcal R\) and generate it as an ordinary group.
Indeed, for a free word in the kernel, insert the representative
of each successive marked-table partial product and its inverse.
The resulting telescoping expression is a product of the displayed
Schreier words; the initial and final representatives are both one.
Normal generation alone would not suffice for the later rank sum.

For every word in (HER.16), enumerate finite products of conjugates
\(v q_i^{\pm1}v^{-1}\) and compare them to that word by free reduction.
Use a fixed dovetail ordering by area and conjugating-word length,
then lexicographic order within each finite stage.  This finds a
certificate for every word: it belongs to the defining normal
closure.  Denote the number of relator factors in its first found
certificate by \(a_w\), and put
\[
       \mathcal D_{\boldsymbol d}=\max\{1,\sum_w a_w\}.
 \tag{HER.17}
\]
Identity words may have area zero.  The finite-table computation
and these finitely many terminating searches specify the promised
algorithm.  They also give verifiable certificates: the group
table shows that all defining relators vanish, while the free-word
certificates show that its ordinary Schreier kernel generators
belong to their normal closure.  Thus the chosen table can be
verified to be the full presented group independently of the
maximum-order search.

Now map \(F\to H\) by the actual labels \(g_1,\ldots,g_k\), and let
\(N\) be the image of \(\mathcal R\).  It is the normal closure in
\(H\) of the power defects \(g_i^{d_i}\).  It is ordinarily generated
by the images \(n_w\) of (HER.16).  From (HER.10), conjugation and
inversion invariance of moving rank, and
\(AB-I=(A-I)B+(B-I)\), each certificate gives
\[
             \operatorname{rank}(n_w-I)\le a_w T.
 \tag{HER.18}
\]
The subspace \(\sum_w\operatorname{im}(n_w-I)\) is exactly
\([N,Q]=\sum_{n\in N}\operatorname{im}(n-I)\): the displayed product
identity bounds each product image by the sum of the two individual
images, and inverse images are the same.  Since \(N\triangleleft H\),
this subspace is \(H\)-invariant.  Irreducibility over the original
\(\mathbf F_p\) implies it is zero or all of \(Q\).
If it is all of \(Q\), (HER.18) yields
\(\dim Q\le\mathcal D_{\boldsymbol d}T\).
If it is zero, the representation factors through
\(\Gamma_{\boldsymbol d}\); the orbit of any nonzero vector spans
the irreducible module, so its dimension is at most \(u\le S\).
There is no semisimplicity or characteristic restriction here.
We have proved
\[
       2m=\dim_{\mathbf F_p}Q
          \le\max\{S,\mathcal D_{\boldsymbol d}T\}
          \le\max\{S,\mathcal D_*(O)T_0\}.
 \tag{HER.19}
\]
Combining this with the low-\(a\) alternative proves the rank bound
in (HER.6).  In particular the argument uses more than the finite
projected signature: the original support gap forces a finite
presentation after low-rank power defects, and its certified
ordinary kernel generators then control the irreducible label space.

\emph{Step 5: projective order and effective finite enumerations.}
The exact sequence supplied by TD05 has kernel \(Q\) and image
\(H\subset\operatorname{Sp}_{2m}(p)\subset\operatorname{GL}_{2m}(p)\).
Consequently
\[
 |P|=p^{2m}|H|
       <p^{(2m)^2+2m}\le B^{K(M)^2+K(M)},
 \qquad a^2=p^{2m}\le B^{K(M)}.
 \tag{HER.20}
\]
Finite group tables below this order bound can be enumerated and
checked.  There is also a direct effective enumeration retaining the
represented Heisenberg action.  For each of the finitely many \(p,m\),
choose a symplectic label basis and the ordinary shift and clock
Weyl matrices.  The nondegenerate twisted group algebra is the full
matrix algebra, as in TD05 and TD12, so this gives its unique
irreducible model up to equivalence for each of the finitely many
nontrivial central characters.  This statement includes \(p=2\).

Its projective normalizer is finite: its kernel on labels is \(Q\)
and its image lies in the finite symplectic group.  To enumerate it
explicitly, enumerate the possible symplectic label maps.  The image
of each order-\(p\) shift or clock generator must be a phase times
the prescribed Weyl matrix, and the order-\(p\) relation leaves only
\(p\) possible phases.  For all these finite choices, solve the
linear intertwining equations.  A nonzero intertwiner is invertible
and unique up to scalar, because both Weyl actions are irreducible.
These exact calculations take place in a cyclotomic number field;
for \(p=2\), the possible fourth-root phases are retained.  They
enumerate every projective normalizer element, including any missing
symplectic label maps by simply finding no solution.  Enumerate all
subgroups of this finite group.  This contains every represented
factor in the theorem up to projective equivalence.  TD63 gives
\(r\le5M\), so all selected marked tuples of the relevant length are
finite enumerations as well.  None of these enumerations is executed
or claimed to decide existence of an original cover or polarization.
\end{proof}

\paragraph{Minimal-normal extraction and the qualified application.}
In the complete TD12 extraction, each Heisenberg factor arises from
a minimal normal elementary abelian subgroup of an actual projective
residual image.  Its nondegenerate represented action gives precisely
(HER.2); the coarse quotient is primitive.  Any nonzero proper
\(H\)-invariant \(\mathbf F_p\)-subspace of \(Q\) would be a nonzero
proper normal subgroup of that chosen minimal normal subgroup.
Thus the extra label-irreducibility hypothesis holds.  All factor
operators and the whole complement are actual quotients of the same
original coefficient action, so (HER.1) is not transferred to a
different source.  The same interfaces hold in the particular accepted
quasiprimitive extraction of TD62.

For the exact minimal, naturally nonsplit, imaginary-quadratic,
rank-nine source class of accepted TD54, its alternating cutoff
gives the complete whole-source floor
\(M_0=\max\{80,N_*-2\}\).  The present theorem therefore gives a
single computable bound for every minimal-normal Heisenberg layer
of that same source.  The numerical definition of \(K\) and its
standalone proof require no alternating-cutoff hypothesis.
Any extension of the application to quasiprimitive sources retains
all exact field, natural determinant and minimal-source qualifications
of the separately reviewed TD62 transfer.

\paragraph{Boundary.}
No finite algorithm is asserted for the remaining geometric existence,
original-field descent, polarization realization or period-map problem.
The new bound concerns this actual selected projective factor under
an explicit whole-original floor and an irreducible original label
module.  It does not replace that hypothesis by a full symplectic
normalizer, and it does not bound original scalar padding or induction.
